\solchapter{optics}

\sol{optics}{1} It appears $2.0/1.33 = 1.5\un{m}$ deep. Light from the whole sky above the water is squeezed into a cone of half-angle $\theta_c = \arcsin(1/1.33) = 49^\circ$ around the vertical (``Snell's window''). Outside that cone the fish sees the surface as a mirror, totally reflecting the underwater scene.

\sol{optics}{2} Diamond: $\arcsin(1/2.42) = 24^\circ$. Glass: $\arcsin(1/1.5) = 42^\circ$. In the fibre, $\sin\theta_c = 1.46/1.48$ gives $\theta_c = 80.6^\circ$ from the normal to the core wall, so light is trapped if it travels within $9.4^\circ$ of the axis. (Equivalently, the fibre accepts light from air within $\arcsin\sqrt{1.48^2-1.46^2} = 14^\circ$ of its axis.)

\sol{optics}{3} $1/s' = 1/10 - 1/30$, so $s' = 15\un{cm}$ and $m = -0.5$: a real, inverted image at half size. At $5\un{cm}$: $1/s' = 1/10-1/5$, so $s' = -10\un{cm}$ and $m = +2$: a virtual, upright, magnified image on the object's side --- a magnifying glass.

\sol{optics}{4} $\Delta y = \lambda L/d = (633\times10^{-9})(2.0)/(10^{-4}) = 1.3\un{cm}$. In water the wavelength shrinks to $\lambda/n = 476\un{nm}$ (the frequency is unchanged), so the fringes close up to $0.95\un{cm}$.

\sol{optics}{5} With $\theta = 1.22\lambda/D$ at $550\un{nm}$: eye ($3\un{mm}$) $2.2\times10^{-4}$, resolving $90\un{km}$ on the Moon; $20\un{cm}$ telescope $3.4\times10^{-6}$, about $1.3\un{km}$; Hubble $2.8\times10^{-7}$, about $100\un{m}$. A $2.4\un{m}$ mirror at $250\un{km}$ resolves $7\un{cm}$ on the ground: enough to see a car, not to read characters whose strokes are about a centimetre wide. The atmosphere's turbulence blurs things further, though looking \emph{down} through it is less damaging than looking up, because the turbulence lies close to the target rather than close to the telescope.

\sol{optics}{6} Reflection off the top surface (air to film) flips the phase; reflection off the bottom (film to air) does not. Constructive interference needs $2nt = (m+\tfrac12)\lambda$, so the thinnest film has $t = \lambda/4n = 550/5.32\approx100\un{nm}$. As $t\to0$, only the half-wavelength flip remains, the two reflections cancel at every wavelength, and the film looks black.

\sol{optics}{7} $\theta_B = \arctan1.33 = 53^\circ$ from the vertical. Glare is then completely horizontally polarized, and is cut most effectively, when the Sun is $90^\circ-53^\circ = 37^\circ$ above the horizon.

\sol{optics}{8} Resolving power is $\lambda/\Delta\lambda = mN$. With $\lambda/\Delta\lambda = 589.3/0.6\approx980$ and $m = 1$, about $1000$ illuminated slits are needed.

\sol{optics}{9} For source $A$ and observer $B$ above a mirror, the path via a point $x$ on the mirror has length $\sqrt{a^2+x^2} + \sqrt{b^2+(d-x)^2}$; minimizing gives $\sin\theta_i = \sin\theta_r$. (Equivalently, reflect $B$ in the mirror; the shortest path to the image is a straight line.) Over a hot road the air is hotter, less dense, and has a slightly lower refractive index, so light travels slightly faster near the road. The least-time path from the sky to your eye dips down toward the road and curves back up, like the lifeguard running along the faster beach. Your eye traces the arriving light back along a straight line and sees the sky on the road: a shimmering ``puddle.''

\sol{optics}{10} Entering, the ray turns by $i - r$; reflecting off the back, by $180^\circ - 2r$; leaving, by $i - r$. Total $D = 180^\circ + 2i - 4r$. Stationary $D$ needs $\dd r/\dd i = \tfrac12$; differentiating Snell's law, $\dd r/\dd i = \cos i/(n\cos r)$, so $n\cos r = 2\cos i$. Combined with $n\sin r = \sin i$, this gives $\cos^2i = (n^2-1)/3$. For $n = 4/3$: $i = 59.4^\circ$, $r = 40.2^\circ$, $D = 138^\circ$, so the bow lies $180^\circ - 138^\circ = 42^\circ$ from the antisolar point. Since $\dd D/\dd i = 0$ there, a wide band of entry angles all exit at nearly the same angle, so light piles up: a caustic.

\sol{optics}{11} Every point on the disc's circular rim is the same distance from the centre of the shadow, so all the secondary wavelets from the rim arrive there in phase. More carefully, divide the unobstructed wavefront outside the disc into Fresnel zones; their contributions alternate in sign with slowly decreasing size, and the sum is about half the first unobstructed zone's contribution --- nearly the same intensity as if no disc were there. Particles would leave a dark shadow. Only waves can make the spot.

\sol{optics}{12} Optics is reversible: any path that light takes from the Sun to the target, light can take from the target back to the Sun. If a passive system could heat the target above the Sun's temperature, the hotter target would then send net heat back to the colder Sun along those same paths --- or, run the other way, heat would flow from cold to hot with no work, violating the second law. The optical statement is conservation of \emph{radiance} (power per area per solid angle): a lens can enlarge the solid angle the Sun fills as seen from the target, but cannot make its surface brighter. The best possible system surrounds the target with the Sun's surface on all sides, which delivers exactly the flux the target would get at the solar surface, and equilibrium is reached at $T_\odot$.
